\documentclass[10pt,a4paper]{amsart}
\usepackage[margin=19mm]{geometry}
\usepackage{amsmath,amssymb,mathtools}
\usepackage[scr=rsfs,cal=euler]{mathalfa}
\usepackage{xfrac}
\usepackage[unicode,hidelinks,bookmarks=false]{hyperref}
\newcommand{\ie}{i.e.\ }
\newcommand{\cf}{cf.\ }
\newcommand{\resp}{respectively\ }
\newcommand{\Subsection}{\subsection}
\providecommand{\nobreakdash}{\nobreak\hbox{-}}
\setlength{\emergencystretch}{2em}
\multlinegap0pt
\newtheorem{lemma}{Lemma}
\newtheorem{prop}{Proposition}
\newcommand\R{\mathbb{R}}
\newcommand\mc{\mathcal}
\title{On the implicit regularization of Langevin dynamics with projected noise}
\author{Govind Menon, Austin J. Stromme, Adrien Vacher}
\date{Source: arXiv:2602.12257v1 (CC BY 4.0)}
\begin{document}
\maketitle

\noindent\emph{Source and licence.} On the implicit regularization of Langevin dynamics with projected noise. Version arXiv:2602.12257v1 (CC BY 4.0), distributed by arXiv under the Creative Commons Attribution 4.0 International licence, http://creativecommons.org/licenses/by/4.0/, as recorded in the arXiv licence metadata for this exact version and shown on the abstract page https://arxiv.org/abs/2602.12257v1. Full licence notice: \texttt{licenses/arxiv-2602.12257-CC-BY-4.0.txt}.\par\smallskip

\noindent\textit{Editorial scope.} Counted unit: the article's lemma lem:second\_fund\_forms\_eq (source0.tex lines 1057--1069). The article's complete original proof of it is delivered with it (source0.tex lines 1531--1584, one \texttt{proof} environment of 54 lines, which the article places away from the statement). No mathematics is written, repaired or re-derived: the statement and the proof are the article's own lines, in the article's own notation, and no retained line is substituted or re-flowed.  Retained uncounted as the source-local context the proof needs: lines 1410--1419.  Nothing was omitted from any retained span, so the package carries no omission record.  No retained line contains a citation, so the wrapper carries no bibliography and no bibliography entry.  No retained line contains a figure environment, an image-inclusion command or drawing code, so no image is delivered and the package declares no asset dependency.  Editorial formatting only, in addition: an AMS \texttt{amsart} preamble that restates the article's own declarations and macros used by the retained lines, namely the environments \texttt{lemma}, \texttt{prop} on separate counters, together with the standard packages those lines need.  The retained environments print in this document as Lemma 1, Proposition 1 in order of appearance, whereas the article prints the same items with the numbering of its own full document.  The complete native source of the article as distributed in the arXiv e-print for this exact version is bundled unchanged as \texttt{sources/194533/source0.tex}.
\begin{prop}[Transformation of horizontal projection]
\label{prop:horizontal_proj}
    Let $P_x$ be the horizontal projection onto the 
    orbit $\mc O_x$ and $Q_x := I - P_x$.
    Then for all $x\in \mathbb{R}^d$ and $g \in G$:
    $$
 gP_x = P_{g\cdot x}g, \quad \quad 
    gQ_x = Q_{g\cdot x} g\,.
    $$
\end{prop}

\begin{lemma}[Relationship between second fundamental forms]
\label{lem:second_fund_forms_eq}
Let $g \in G$ and $x \in \R^d$.
Let $h_{g \cdot x}^{\mc O_x}$ be the second
fundamental form of $\mc O_x$ at $g \cdot x$
and $h^G_g$ be the second fundamental form
of $G$ at $g$.
Then for all $A \in T_{g} G$ we have
$$
h_{g \cdot x}^{\mc O_x} (L_{g, x}A,L_{g,x}A) = 
P_{g \cdot x}(h^G_g(A, A) x).
$$
\end{lemma}

\begin{proof}[Proof of Lemma~\ref{lem:second_fund_forms_eq}]
We first compute the second fundamental
form $h^G_g(A, A)$, for $A \in T_gG$. Let us begin
by considering the case where $g = I$. Given
any smooth vector field $W$ on $\R^d$ with
$W_I = A$ and such that $W_g \in T_gG$ for all $g\in G$, the second fundamental form is
$$
h^G_I(A, A) = P^G_I(\nabla_{W} W),
$$ where $P^G_I$ is the projection form $T_I\R^d$ to the normal
space $N_IG$ and $\nabla$ denotes the Euclidean connection.
We may therefore take $W$ to be $W_X:= XA$ for all $X \in \mathbb{R}^{d \times d}$,
which will be a 
smooth vector field on $G$ and such that $W_I = A$. Then
$$
(\nabla_{W}W)_I = \lim_{t\to 0} \frac{W_{I + tA} - W_{I}}{t}
= A^2.
$$ But then notice that, since $G \subset O(d)$, $A$ must be anti-symmetric,
so that in fact the above is a symmetric matrix. Since
$T_gG$ is a subset of anti-symmetric matrices, it follows that
$\nabla_W W \in N_IG$ already. Hence
$$
h^G_I(A, A) = A^2.
$$ We may then deduce the general formula by the fact that
$g \colon \R^{d \times d} \to \R^{d \times d}$ which takes
$X \mapsto gX$, is an isometry so that for $B \in T_gG$, we have
$$
h^G_g(B, B) = gh^G_I(g^{-1}B, g^{-1}B) = Bg^{-1}B.
$$

Next, we compute the second fundamental form of $\mc O_x$.
Here, we know that each $v \in T_x \mc O_x$ is of the form
$Ax$ for $A \in T_eG$. We can then extend $v$ to the vector
field $V(x') := Ax'$, and again obtain
$$
(\nabla_V V)_x = \lim_{t\to 0} \frac{V(x + tAx) - V(x)}{t}
= A^2x.
$$ Hence if $v = Ax$ then
\begin{align*}
h^{\mc O_x}_x(v, v)= P_x(\nabla_V V)_x = P_x(A^2x) = P_x(h^G_I(A, A)x)&=
P_x(g^{-1}h^G_g(gA, gA)x),
\end{align*}
yielding the result.
Applying Proposition~\ref{prop:horizontal_proj}, this implies
$$
gh_x^{\mc O_x}(v, v) = P_{g\cdot x}
(h_g^G(gA, gA)x).
$$ Since $g \colon \R^d \to \R^d$ is an isometry the left-hand-side
becomes
$$
h_{g\cdot x}^{\mc O_x}(g\cdot v, g \cdot v)
= P_{g\cdot x}(h_g^G(gA, gA)x).
$$ Because $v$ was an arbitrary element of $T_x \mc O_x$,
the result follows.
\end{proof}

\end{document}
